%% =====================================================================
%%  B4-T-21-04 -- what B4 contributes to "The Frozen Crowd"
%%  -------------------------------------------------------------------
%%  LAYER A: APPEND-ONLY. Written once at the entry's write, then left
%%  alone. If the source has more to say later, add a bullet -- do not
%%  rewrite. Material found ONLY here goes in the `unique` slot with
%%  @unique set to yes, which makes it undeletable (rule R10).
%% =====================================================================
%% @kind: source
%% @id: B4-T-21-04
%% @book: B4
%% @topic: T-21-04
%% @locator: ch. 4 (Social Proof), pp. 126--156 --- the bystander section, pp. 97--110
%% @status: distilled
%% @unique: yes
%% @integrated: yes
%% @updated: 2026-09-19

\dossier{B4}{T-21-04}{\bookshort{B4}}{ch. 4 (Social Proof), pp. 126--156 --- the bystander section, pp. 97--110}

\begin{distilled}
  \item Pluralistic ignorance: in ambiguity everyone checks everyone else, and each person's concealed alarm is misread as calm --- thirty-eight neighbours watched the Genovese murder through lit windows and none called in time.
  \item Diffusion of responsibility: the more witnesses, the less weight on any single one; help addressed to a group is addressed to nobody.
  \item The laboratory dose-response curve: victims get help fastest one-on-one, slowest in crowds.
  \item The countermeasure tested in the field: individuate --- a pointed request with a name and eye contact dissolves the diffusion instantly.
  \item B4's structure: the same social-proof machinery that spreads panic spreads inaction; the crowd's uncertainty is the mechanism's fuel.
\end{distilled}
\begin{unique}
  \item The Genovese case, pluralistic ignorance as a named mechanism, and the dose-response bystander experiments are unique to B4 in this corpus.
\end{unique}
\begin{terms}
  \item \emph{pluralistic ignorance} --- B4's term for a crowd each of whose members privately doubts while reading everyone else's calm face as reassurance.
  \item \emph{diffusion of responsibility} --- the dilution of any individual's sense of duty in proportion to the number of witnesses.
\end{terms}
